\documentclass[11pt]{article}
\usepackage[margin=1in]{geometry}
\usepackage[T1]{fontenc}
\usepackage[utf8]{inputenc}
\usepackage{amsmath,amssymb,amsthm,array,booktabs,microtype,enumitem}
\usepackage{cite}\renewcommand{\citemid}{,~}
\usepackage[hidelinks]{hyperref}
\hypersetup{pdftitle={The Cranston\textendash Li\textendash Su\textendash Wang\textendash Xu (2026) conjecture is true},pdfauthor={Dave Goel},
  pdfkeywords={nowhere-zero flow, group flow, reconfiguration, flow-connected graph, Eulerian graph}}

\newtheorem{theorem}{Theorem}[section]
\newtheorem{lemma}[theorem]{Lemma}
\newtheorem{corollary}[theorem]{Corollary}
\newtheorem{proposition}[theorem]{Proposition}
\theoremstyle{definition}
\newtheorem{definition}[theorem]{Definition}
\newtheorem{remark}[theorem]{Remark}
\newcommand{\Z}{\mathbb{Z}}
\newcommand{\cF}{\mathcal{F}}
\newcommand{\cG}{\mathcal{G}}
\newcommand{\one}{\mathbf{1}}
\raggedbottom
\makeatletter\@secpenalty=0\makeatother
\widowpenalty=10000 \clubpenalty=10000
\emergencystretch=1.5em

\title{The Cranston--Li--Su--Wang--Xu (2026) conjecture is true}
\author{Dave Goel\thanks{\raggedright The Lean proof is at
\url{https://github.com/david44hello/papers_supporting_materials/tree/cranston-li-su-wang-xu-2026-conjecture-v1.59/cranston-li-su-wang-xu-2026-conjecture}.}}
\date{October 2026}

\begin{document}
\maketitle

\begin{abstract}
Let $A$ be a finite abelian group. A nowhere-zero $A$-flow of a graph gives every edge a direction and a non-zero
value in $A$, so that at every vertex the sum of the values on the entering edges equals the sum on the leaving edges.
Such a flow can be changed along a cycle: going around the cycle, add a fixed non-zero element of $A$ to the value of
every edge that points forward and subtract it from the value of every edge that points backward. The Cranston--Li--Su--Wang--Xu (2026) conjecture states that if $A$ has at least four elements, then in every
graph in which all degrees are even, any two nowhere-zero
$A$-flows are joined by a chain of such changes through nowhere-zero $A$-flows. It also states that, for every integer $k\ge 4$, the same holds for integer flows whose
values are between $1$ and $k-1$ in absolute value, where each change adds a non-zero integer. We prove this conjecture. The core of the proof is an induction on the number of edges: if a flow
enters a vertex $v$ along an edge $uv$ and leaves $v$ along an edge $vw$ with the same value, the two edges can be
replaced by one edge $uw$, and the induction hypothesis applies to the smaller graph. So the main task is to change
every flow into a flow with such a pair of edges.
For nowhere-zero $A$-flows with $|A|\ge 4$, replacing vertices of degree more than $4$ by triangles reduces the problem to graphs with all degrees
at most $4$. For integer flows, a counting argument joins any two nowhere-zero flows with values in the cyclic group $\Z_k$,
$k\ge 4$, by changes that add $1$, $-1$, $2$ or $-2$; every such chain gives a chain of integer flows with values between
$1$ and $k-1$ in absolute value, and every integer flow of this chain reduces modulo $k$ to a flow of the given chain. All results of this paper are also proved in the Lean~4 proof
assistant, from Lean's standard axioms alone.
\end{abstract}

\noindent\textit{2020 Mathematics Subject Classification:} 05C21 (primary), 05C45 (secondary).\\
\textit{Keywords:} nowhere-zero flow, group flow, reconfiguration, flow-connected graph, Eulerian graph.

\section{Introduction}\label{sec:intro}

This paper proves the Cranston--Li--Su--Wang--Xu (2026) conjecture \cite[Conjecture~4.1]{CLSWX} on the
reconfiguration of nowhere-zero flows in graphs in which every vertex has even degree, that is, on whether every
nowhere-zero flow can be changed into every other one by changes along cycles that keep the flow nowhere-zero.

Let $G$ be a finite graph; loops and parallel edges are allowed. For a finite abelian group $A$ (for example the cyclic group $\Z_m$ of the integers modulo $m$), a nowhere-zero
$A$-flow of $G$ assigns to every edge, in a fixed reference direction, a non-zero element of $A$, so that at every
vertex the sum of the values on the entering edges equals the sum on the leaving edges. Two such flows differ by a
\emph{step} with \emph{value} $t\ne 0$ if there is a cycle, traversed in one of its two directions, such that their
difference is $t$ on every edge of the cycle whose reference direction agrees with the traversal, $-t$ on every other
edge of the cycle, and $0$ on every edge not on the cycle. The graph $\cF(G,A)$ has the nowhere-zero $A$-flows of $G$ as its vertices, two of them being
adjacent if they differ by a step, and $G$ is \emph{$A$-flow-connected} if $\cF(G,A)$ is connected. For an integer
$k\ge 2$, the graph $\cF(G,k)$ is defined in the same way with integers. Its vertices are the integer flows $F$ of $G$
(flows with integer values) with $1\le|F(e)|\le k-1$ on every edge $e$, and two of them are adjacent if they differ by a
step whose value is a non-zero integer. The graph $G$ is \emph{$k$-flow-connected} if $\cF(G,k)$ is connected. Sections~\ref{sec:words} and~\ref{sec:int}
restate these definitions in the notation of the proofs.

Esperet, Hendrey, Lagoutte, Marseloo, Norin and Steiner \cite{EHLMNS} introduced the graphs $\cF(G,A)$ and $\cF(G,k)$
and asked for which graphs $G$, groups $A$ and integers $k$ they are connected. Cranston, Li, Su, Wang and Xu
\cite[Conjecture~4.1]{CLSWX} conjectured that every graph in which all degrees are even is $A$-flow-connected for every finite abelian group $A$
with $|A|\ge 4$ and $k$-flow-connected for every integer $k\ge 4$. In \cite{CLSWX}, a graph in which all
degrees are even is called Eulerian. Since ``Eulerian'' often also requires the graph to be connected, we call a graph \emph{even} if every vertex
has even degree; an even graph need not be connected.

\begin{theorem}\label{thm:main}
Let $G$ be a graph in which every vertex has even degree.
\begin{itemize}
\item[(a)] $G$ is $A$-flow-connected for every finite abelian group $A$ with $|A|\ge 4$.
\item[(b)] $G$ is $k$-flow-connected for every integer $k\ge 4$.
\end{itemize}
\end{theorem}

Thus the Cranston--Li--Su--Wang--Xu (2026) conjecture is true. Part~(a) is Theorem~\ref{thm:all}, proved in
Section~\ref{sec:proofs}, and part~(b) is Theorem~\ref{thm:int}, proved in Section~\ref{sec:int}. The following table
compares Theorem~\ref{thm:main} with earlier results.

\begin{center}
\small
\begin{tabular}{@{}>{\raggedright\arraybackslash}p{0.45\textwidth}@{\hspace{0.04\textwidth}}>{\raggedright\arraybackslash}p{0.51\textwidth}@{}}
\toprule
New result & Best result before this paper\\
\midrule
Every even graph is $A$-flow-connected for every finite abelian group $A$ with $|A|\ge 4$
(Theorem~\ref{thm:main}~(a)).
& Groups of even order \cite[Theorem~4.3]{CLSWX}; groups with $|A|\ge 5{,}251{,}339$
\cite[Corollary~4.5]{CLSWX}; planar graphs, $|A|\ge 7$ \cite[Corollary~5.9]{EHLMNS}.\\
\addlinespace
Every even graph is $k$-flow-connected for every integer $k\ge 4$ (Theorem~\ref{thm:main}~(b)).
& Even $k$ \cite[Theorem~4.2]{CLSWX}; planar graphs, $k\ge 7$ \cite[Corollary~5.9]{EHLMNS}; open for odd $k\ge 5$ in general.\\
\bottomrule
\end{tabular}
\end{center}

The number $4$ in Theorem~\ref{thm:main} cannot be replaced by $3$. By \cite[Theorem~2.1]{CLSWX}, for a graph with a
nowhere-zero $3$-flow (an integer flow with values in $\{1,-1,2,-2\}$), the following are equivalent: it is
$\Z_3$-flow-connected; it is $3$-flow-connected; all its degrees are even and no two of its vertices are joined by more
than four edge-disjoint paths. The even graph with two vertices and six parallel edges has a nowhere-zero $3$-flow
(direct every edge from one vertex to the other, and give three edges the value $1$ and three the value $-1$), and its
two vertices are joined by six edge-disjoint paths. So it is neither $\Z_3$-flow-connected nor $3$-flow-connected.

\paragraph{Further results.} Every even graph in which all degrees are at most $4$ is $A$-flow-connected for every
finite abelian group $A$, of any order (Theorem~\ref{thm:deg4}); for $A=\Z_3$ this also follows from
\cite[Theorem~2.1]{CLSWX}, since an even graph has a nowhere-zero $3$-flow (the value $1$ on every edge, with the edges directed so
that at every vertex as many edges enter as leave, which is possible because all degrees are even) and, when all degrees are at most $4$, no two
vertices are joined by more than four edge-disjoint paths. For the cyclic groups $\Z_m$ with $m\ge 4$, any two nowhere-zero $\Z_m$-flows of an even graph are joined by a
chain of steps whose values are in $\{1,-1,2,-2\}$ (Theorem~\ref{thm:cyclic}); for even $m$, $\Z_m$-flow-connectedness, with steps of any value,
was known \cite[Theorem~4.3]{CLSWX}. The restriction of the step values to $\pm1$ and $\pm2$ is what the proof of
part~(b) of Theorem~\ref{thm:main} needs (see the outline below).

An expanded version of this paper, with many other results, is posted on Zenodo \cite{G}.
The two main questions of combinatorial reconfiguration, as listed in \cite[Section~1]{EHLMNS}, are whether the
reconfiguration graph (here $\cF(G,A)$) is connected and, if so, what its diameter is. For even graphs and $|A|\ge 4$,
Theorem~\ref{thm:main} answers the first. For the second, \cite[Theorem~2.5]{G} gives a polynomial upper bound: if $|A|\ge 4$, any two
nowhere-zero $A$-flows of an even graph with $M$ edges are joined by a chain of at most $16M^3+12M^2+M$ steps. For
integers $q\ge 1$ and $p>2q$, a modular $(p,q)$-flow \cite[Section~2]{CLSWX} is a $\Z_p$-flow whose values are residues
of $q,q+1,\dots,p-q$, and a graph is $(p,q)$-flow-connected if any two of its modular $(p,q)$-flows are joined by a
chain of steps through modular $(p,q)$-flows. All even graphs are $(p,q)$-flow-connected if and only if $p\ge 4q$
\cite[Theorems~9.1~and~9.2]{G}. None of these results is used here.

\paragraph{Outline of the proof.}
Section~\ref{sec:red} reduces the flow-connectedness of all even graphs to the statement that, in every loopless even graph
with at least one edge, every flow can be changed, by a chain of steps, into a flow with a balanced pair (defined next) at some vertex. If a flow enters a vertex $v$ along an edge
$e=uv$ and leaves along another edge $e'=vw$ with the same value (a balanced pair, Definition~\ref{def:balanced}), the two
edges can be replaced by a single edge $uw$, and the flows with this balanced pair correspond to the flows of the smaller
graph (Lemma~\ref{lem:split}). Assuming the statement above, we use induction on the number of edges: the smaller graph is flow-connected by the
induction hypothesis, so any two flows with a balanced pair on the same two edges are joined by a chain of steps. An Eulerian orientation directs every edge so that at every
vertex as many edges enter as leave; it exists because every degree is even. A constant flow with value $a\ne 0$ has
the value $a$ on every edge in the direction of an Eulerian orientation. For any two edges at a vertex of a loopless even graph and any $a\ne 0$,
the two edges form a balanced pair at that vertex in some constant flow with value $a$, and for a suitable $a$, all constant flows with
value $a$ are joined to each other (Lemma~\ref{lem:const}). So all flows with a balanced pair are joined to each other, and it
suffices to join every flow to one of them; this is statement \textup{(LOCAL)} of Proposition~\ref{prop:red}.

Statement \textup{(LOCAL)} is needed only for loopless graphs in which every degree is $0$ or at least $4$: loops can be
deleted (Lemma~\ref{lem:loop}), and at a vertex of degree $2$ the two edges form a balanced pair in every flow, by the flow condition.
Section~\ref{sec:local} proves it when every degree is $0$ or $4$; this gives Theorem~\ref{thm:deg4}
(all degrees at most~$4$). In Section~\ref{sec:proofs} a vertex of larger degree is replaced by a triangle: if
$|A|\ge 4$ and the graph with the triangle is $A$-flow-connected, then so is the original graph
(Lemma~\ref{lem:triangle}). This implication is a special case of \cite[Theorem~5.9(1)]{CLSWX}; we include its short
proof. That theorem applies because a triangle is $A$-connected when $|A|\ge 4$, in the sense of \cite{JLPT}. This
notion is different from $A$-flow-connectedness: a graph $H$ is \emph{$A$-connected} if, for a fixed reference
direction of every edge and every map $\beta\colon V(H)\to A$ with $\sum_v\beta(v)=0$, there is a map
$f\colon E(H)\to A\setminus\{0\}$ such that at every vertex $v$ the sum of $f$ on the leaving edges minus the sum on the
entering edges is $\beta(v)$.

Part~(b) of
Theorem~\ref{thm:main} needs a further argument: the steps of Section~\ref{sec:local} can have any value, but
Lemmas~\ref{lem:step1} and~\ref{lem:step2}, which turn a step between two $\Z_k$-flows into a step between two integer flows
whose residues modulo $k$ are these two $\Z_k$-flows, apply only to steps with values $\pm1$
and $\pm2$. So, for the cyclic groups $\Z_m$ with $m\ge 4$,
Section~\ref{sec:cyclic} proves \textup{(LOCAL)} using only steps with values $\pm1$ and $\pm2$, by a counting argument
that needs no bound on the degrees; this gives Theorem~\ref{thm:cyclic}, which also proves part~(a) for cyclic groups
(Remark~\ref{rem:noncyclic} explains why the counting argument does not extend to all groups). Section~\ref{sec:int} lifts these chains to
integer flows, using Tutte's lifting theorem \cite{Tutte} in the form of \cite[Theorem~2.3]{EHLMNS}: every nowhere-zero
$\Z_k$-flow is the residue modulo $k$ of an integer flow with values between $1$ and $k-1$ in absolute value.

\section{Definitions and results}\label{sec:words}

\paragraph{Graphs.} All graphs are finite; loops and parallel edges are allowed. A graph is \emph{even} if every
vertex has even degree, where a loop counts twice. A \emph{cycle} is a connected subgraph in which every vertex has
degree~$2$; a loop is a cycle with one edge, and two parallel edges form a cycle with two edges. A \emph{directed
cycle} is a cycle together with one of its two directions of traversal.

\paragraph{Flows.} Let $A$ be a finite abelian group, and let $G$ be a graph. Fix a
reference direction of every edge. An \emph{$A$-flow} of $G$ is a map $f\colon E(G)\to A$ such that at every vertex
the sum of $f$ over the edges that leave the vertex equals the sum over the edges that enter it. The value of $f$ on an
edge $e$ is $f(e)$ in the reference direction of $e$ and $-f(e)$ in the other direction. From now on, a \emph{flow}
is a nowhere-zero $A$-flow, that is, $f(e)\ne 0$ for every edge $e$.

For a directed cycle $Q$, let $\chi_Q(e)$ be $1$ if $Q$ traverses $e$ in its reference direction, $-1$ if
$Q$ traverses $e$ against it, and $0$ if $e$ is not on $Q$. A \emph{step} is a pair of flows $f$ and
$f'=f+t\chi_Q$, where $Q$ is a directed cycle and $t\ne 0$ is an element of $A$, the \emph{value} of the step. The value is determined only up to sign, as reversing $Q$ or exchanging $f$ and $f'$ replaces $t$ by $-t$. The
graph $\cF(G,A)$ has the flows of $G$ as its vertices, and two flows are adjacent if they form a step. In \cite{CLSWX,EHLMNS} two different flows are adjacent if they agree outside a cycle $Q$; this gives the same graph, because the
difference of two such flows is an $A$-flow that is zero outside $Q$, and such an $A$-flow is $t\chi_Q$ for some $t\in A$ and a direction of $Q$. The graph $G$ is \emph{$A$-flow-connected} if $\cF(G,A)$ is connected; here $\cF(G,A)$ counts as connected when it has at most one
vertex, in particular when $G$ has no~flow.

Let $T$ be a set of non-zero elements of $A$ with $T=-T$. A \emph{$T$-step} is a step whose value is in $T$, and two flows
are \emph{$T$-joined} if a chain of $T$-steps (possibly empty) leads from one to the other.

A \emph{multi-step} with value $t\ne 0$ is a pair of flows $f$ and $f'=f+t(\chi_{Q_1}+\dots+\chi_{Q_r})$, where
$r\ge 1$ and $Q_1,\dots,Q_r$ are directed cycles, no two of which have a common edge. Every partial sum
$f+t(\chi_{Q_1}+\dots+\chi_{Q_j})$ is then a flow, as it agrees with $f$ or $f'$ on every edge; so a
multi-step is a chain of $r$ steps with value $t$.

\paragraph{Ends of edges.} Every edge has two ends, one at each endpoint; a loop at $v$ has both ends at $v$. The degree of $v$ is the
number of ends at $v$. If $v$ is a vertex of a cycle $Q$, the \emph{ends of $Q$ at $v$} are the two ends at $v$ of the
edges of $Q$. For an $A$-flow $f$ and an end $i$ at $v$ of the edge $e$, the \emph{outward value}
$y_i(f)\in A$ is the value of $f$ on $e$ in the direction that leaves $v$ through $i$; it is non-zero if $f$ is a flow. The flow
condition at $v$ says
\begin{equation}\label{eq:kirchhoff}
\sum_{i\text{ an end at }v} y_i(f)=0 .
\end{equation}
If a directed cycle $Q$ enters $v$ through the end $i$ and leaves $v$ through the end $j$, then
\begin{equation}\label{eq:ends}
y_i(f+t\chi_Q)=y_i(f)-t,\qquad y_j(f+t\chi_Q)=y_j(f)+t ,
\end{equation}
and the other ends at $v$ keep their outward values.

\begin{definition}\label{def:balanced}
Let $G$ have no loop. A \emph{balanced pair} of a flow $f$ at a vertex $v$ is a pair of different ends $i,j$ at
$v$ with $y_i(f)+y_j(f)=0$; that is, $f$ enters $v$ through one of them and leaves $v$ through the other with the same
value. (Loops are excluded because the two ends of a loop would always form a balanced pair.)
\end{definition}

We prove the following three theorems on group flows, that is, on $A$-flows for a finite abelian group $A$.
\begin{theorem}\label{thm:deg4}
Let $A$ be a finite abelian group. Every even graph in which all degrees are at most $4$ is $A$-flow-connected.
\end{theorem}

\begin{theorem}\label{thm:all}
Let $A$ be a finite abelian group with $|A|\ge 4$. Every even graph is $A$-flow-connected.
\end{theorem}

\begin{theorem}\label{thm:cyclic}
Let $m\ge 4$. In every even graph, any two nowhere-zero $\Z_m$-flows are joined by a chain of steps whose values are
in $\{1,-1,2,-2\}$.
\end{theorem}

Sections~\ref{sec:proofs} and~\ref{sec:cyclic} prove these theorems, and Section~\ref{sec:int} deduces Theorem~\ref{thm:main}~(b) from Theorem~\ref{thm:cyclic}.

\section{Reduction to flows with a balanced pair}\label{sec:red}

This section proves the following (Proposition~\ref{prop:red}). Let a class of even graphs be closed under deleting a
loop and under splitting off, that is, replacing two edges $uv$ and $vw$ at a vertex $v$ by one edge $uw$
(Lemma~\ref{lem:split}). If, in every loopless graph of the class with at least one edge, every flow is joined by steps
to a flow with a balanced pair, then in every graph of the class any two flows are joined by steps. In this section $A$ is a finite abelian group with at least two elements, and $T$ is a set of non-zero elements of $A$
with $T=-T$. Here $T$ is the set of allowed step values ($A\setminus\{0\}$ for Theorems~\ref{thm:deg4} and~\ref{thm:all}, and
$\{1,-1,2,-2\}$ for Theorem~\ref{thm:cyclic}). We use two conditions on $T$:
\begin{itemize}
\item[(T1)] any two non-zero elements of $A$ are joined by a sequence of non-zero elements of $A$ in which
consecutive elements differ by an element of $T$;
\item[(T2)] there is a non-zero $a\in A$ with $2a=0$ or $2a\in T$.
\end{itemize}
Both conditions hold for $T=A\setminus\{0\}$, and also for $A=\Z_m$ with $m\ge 3$ and $T=\{1,-1,2,-2\}$: the sequence
$1,2,\dots,m-1$ gives (T1), and $a=1$ gives (T2). Condition (T1) is used for loops in Lemma~\ref{lem:loop}, and (T2) is
used in Proposition~\ref{prop:red} to join the constant flows of Lemma~\ref{lem:const}.

\begin{lemma}[deleting a loop]\label{lem:loop}
Let $T$ satisfy \textup{(T1)}, and let $\ell$ be a loop of $G$. If any two flows of $G-\ell$ are $T$-joined, then any two
flows of $G$ are $T$-joined.
\end{lemma}
\begin{proof}
A loop adds the same term to both sides of the flow condition at its vertex. So a flow of $G$ is a flow of $G-\ell$
together with any non-zero value on $\ell$. A step of $G-\ell$, with the value on $\ell$ kept, is a step of $G$, and a
change of the value on $\ell$ by an element of $T$ is a $T$-step along the cycle $\ell$. Given two flows of $G$, join their
restrictions by $T$-steps of $G-\ell$, then change the value on $\ell$ along a sequence as in \textup{(T1)}.
\end{proof}

\begin{lemma}[splitting off]\label{lem:split}
Let $G$ have no loop, let $v$ be a vertex, and let $e=vu$ and $e'=vw$ be two different edges at $v$ ($u=w$ is
allowed). Let $G'$ be the graph $G-e-e'+g$ with a new edge $g=uw$ (a loop if $u=w$): the first end of $g$ is at $u$ and
replaces the end of $e$ at $u$, and its second end is at $w$ and replaces the end of $e'$ at $w$. The direction of
$g$ \emph{from $u$ to $w$} is the direction from its first end to its second end; if $u=w$, this fixes one of the two
directions of the loop $g$. For a map $f'\colon E(G')\to A$, let $x$ be the value of $f'$ on $g$
from $u$ to $w$, and let $Lf'\colon E(G)\to A$ be the map that agrees with $f'$ on the edges other than $e,e',g$ (the common edges of $G$ and $G'$) and has the value $x$ on
$e$ in the direction from $u$ to $v$ and on $e'$ in the direction from $v$ to $w$. Let $B(v;e,e')$ be the set of flows
$f$ of $G$ in which the ends of $e$ and $e'$ at $v$ are a balanced pair.
\begin{itemize}
\item[(a)] $L$ is a bijection from the flows of $G'$ onto $B(v;e,e')$.
\item[(b)] If $f'_2$ comes from $f'_1$ by a multi-step of $G'$ with value $t$, then $Lf'_2$ comes from $Lf'_1$ by
a multi-step of $G$ with value $t$. In particular, if two flows of $G'$ are $T$-joined, their images under $L$ are
$T$-joined.
\end{itemize}
\end{lemma}
\begin{proof}
(a) In $Lf'$ the ends of $e$ and $e'$ at $v$ have the outward values $-x$ and $x$, so the flow condition holds at $v$;
at $u$ the edge $e$ leaves with value $x$ as $g$ did, and at $w$ the edge $e'$ enters with value $x$ as $g$
did. So $Lf'$ is a flow in $B(v;e,e')$. Conversely, for $f\in B(v;e,e')$ the value of $f$ on $e$ from $u$ to $v$
equals its value on $e'$ from $v$ to $w$; giving $g$ this value from $u$ to $w$ and keeping the other values defines
the only map $f'$ with $Lf'=f$; it is a flow, as the ends of $g$ have the outward values of the ends of $e,e'$ that they
replace, and the ends of $e,e'$ at $v$ have outward values with sum~$0$.

(b) Let $Q'$ be a directed cycle of $G'$. If $g$ is not on $Q'$, let $W=Q'$, a directed cycle of $G$. If $g$ is on
$Q'$, replace $g$ by the two edges $e,e'$, traversed through $v$ in the direction of $Q'$ (first $e$ from $u$ to $v$
and then $e'$ from $v$ to $w$ if $Q'$ traverses $g$ from $u$ to $w$, and the reverse otherwise); this gives a closed
walk $W$ of $G$ that uses every edge at most once. If $u=w$, then $Q'=g$ and $W$ is the cycle formed by $e$ and $e'$.
If $u\ne w$ and $Q'$ does not pass through $v$, then $W$ is a directed cycle. If $u\ne w$ and $Q'$ passes through $v$,
then $W$ visits $v$ twice and no other vertex twice, and splitting $W$ at $v$ gives two directed cycles with no
common edge. In every case let $\chi_W$ be the sum of the maps $\chi$ of these one or two directed cycles; then
$L(f'+t\chi_{Q'})=Lf'+t\chi_W$. If $Q'_1,\dots,Q'_r$ are directed cycles of $G'$, no two of which have a common edge,
then the directed cycles of $G$ obtained from them have pairwise no common edge: they use the edges of the $Q'_j$,
with $e,e'$ instead of $g$. So $L$ maps a multi-step of $G'$ to a multi-step of $G$ with the same value.
\end{proof}

\begin{lemma}[constant flows]\label{lem:const}
Let $G$ be even and let $a\ne 0$ be an element of $A$. An Eulerian orientation $D$ of $G$ directs every edge so that at every vertex the
in-degree equals the out-degree (one exists because $G$ is even); the directed edges are the \emph{arcs} of $D$. Let
$a\one_D$ be the constant flow with value $a$ along every arc of $D$.
\begin{itemize}
\item[(a)] For two Eulerian orientations $D,D_2$ of $G$, either $a\one_{D_2}=a\one_D$, or $a\one_{D_2}$ comes from
$a\one_D$ by a multi-step with value $-2a$.
\item[(b)] In the situation of Lemma~\ref{lem:split}, $B(v;e,e')$ contains $a\one_D$ for some Eulerian orientation $D$ of $G$.
\end{itemize}
\end{lemma}
\begin{proof}
(a) The arcs of $D$ that are reversed in $D_2$ form a directed graph with in-degree equal to out-degree at every
vertex, so they are the union of directed cycles $Q_1,\dots,Q_r$ of $D$, no two of which have a common arc. The
reversal of an arc changes its value along $D$ from $a$ to $-a$. So
$a\one_{D_2}=a\one_D-2a(\chi_{Q_1}+\dots+\chi_{Q_r})$. If $2a=0$ or $r=0$, the two flows are equal; otherwise this is a
multi-step with value $-2a$.

(b) The graph $G'$ is even. Take an Eulerian orientation $D'$ of $G'$, and let $D$ agree with $D'$ on the common edges and
direct $e$ and $e'$ through $v$ in the direction of $g$ in $D'$ (from $u$ to $v$ to $w$ if $D'$ directs $g$ from $u$
to $w$, and the reverse otherwise). Then $D$ is an Eulerian orientation of $G$ and $a\one_D=L(a\one_{D'})$.
\end{proof}

\begin{proposition}[reduction to a balanced pair]\label{prop:red}
Let $\cG$ be a class of even graphs such that $G-\ell\in\cG$ for every $G\in\cG$ and every loop $\ell$ of $G$, and
$G'\in\cG$ for every $G\in\cG$ with no loop and every graph $G'$ obtained from $G$ as in Lemma~\ref{lem:split}. Let
$T$ satisfy \textup{(T1)} and \textup{(T2)}. Suppose that the following statement holds.
\begin{itemize}[leftmargin=4.6em,labelsep=0.5em]
\item[\textup{(LOCAL)}] For every loopless $G\in\cG$ in which every vertex has degree $0$
or at least $4$, and which has at least one edge, every flow of $G$ is $T$-joined to a flow that has a balanced pair at some vertex.
\end{itemize}
Then in every $G\in\cG$ any two flows are $T$-joined.
\end{proposition}
\begin{proof}
We use induction on the number of edges of $G$. A graph with no edge has exactly one flow. If $G$ has a loop $\ell$, use
Lemma~\ref{lem:loop} and the induction hypothesis for $G-\ell$. From now on $G$ has no loop.

If some vertex $v$ has degree $2$, its two ends belong to two different edges $e,e'$, and by \eqref{eq:kirchhoff}
they are a balanced pair in every flow. So every flow of $G$ is in $B(v;e,e')$, and by Lemma~\ref{lem:split} and the
induction hypothesis for $G'$, any two flows of $G$ are $T$-joined.

Otherwise $G$ has an edge and every vertex has degree $0$ or at least $4$. Fix $a$ as in \textup{(T2)}. For every
vertex $v$ and any two different edges $e,e'$ at $v$, any two flows in $B(v;e,e')$ are $T$-joined, by the induction
hypothesis for $G'$ and Lemma~\ref{lem:split}, and $B(v;e,e')$ contains a flow $a\one_D$, by
Lemma~\ref{lem:const}~(b). By Lemma~\ref{lem:const}~(a) and \textup{(T2)}, all flows $a\one_D$ are $T$-joined to one another (if $2a\ne 0$, the value
$-2a$ is in $T$, because $2a\in T$ and $T=-T$). A flow with a balanced pair is in some $B(v;e,e')$, because $G$ has
no loop. So all flows with a balanced pair are $T$-joined to one another, and \textup{(LOCAL)} joins every flow to
one of them.
\end{proof}

\section{Reaching a balanced pair when all degrees are 0 or 4}\label{sec:local}

This section proves statement \textup{(LOCAL)} of Proposition~\ref{prop:red} for graphs in which every
degree is $0$ or $4$ (Proposition~\ref{prop:local4}). We fix generators of the cyclic summands of $A$ and call an edge
small if its value is one of these generators or its negative (precise definitions follow). First, steps along cycles
change a flow into one in which the edges that are not small form a forest (Corollary~\ref{cor:forest}); then, as the graph has an edge and every vertex with an edge has degree $4$, the small
edges contain a cycle, and one more step along that cycle gives a balanced pair (Lemma~\ref{lem:Bstar}). Let $A$ be a finite abelian
group with at least two elements.
Write $A$ as a direct sum of cyclic groups,
\[
A=\langle s_1\rangle\oplus\dots\oplus\langle s_r\rangle,\qquad s_h\ne 0 ,
\]
so that every element $a$ of $A$ can be written in exactly one way as $a=a_1+\dots+a_r$ with $a_h\in\langle
s_h\rangle$; the $a_h$ are the \emph{parts} of $a$. Let $S=\{s_1,-s_1,\dots,s_r,-s_r\}$. The elements of $S$ are
called \emph{small}, and the \emph{color} of $s\in S$ is the index $h$ with $s\in\langle s_h\rangle$. We have $S=-S$ and $0\notin S$, and $S$ generates $A$. For a
flow $f$, an edge is \emph{small} if its value is small (this does not depend on the direction); the \emph{color} of
a small edge is the color of its value, which is the same in both directions. Let $\rho(f)$ be the number of edges
that are not small.

\begin{lemma}[a cycle with no small edge]\label{lem:Astar}
Let $f$ be a flow of a graph $G$ and $Q$ a directed cycle of $G$ none of whose edges is small. Then there is
a step $f'=f+t\chi_Q$ with $\rho(f')\le\rho(f)-1$.
\end{lemma}
\begin{proof}
Let $X$ be the set of values of the edges of $Q$ in the direction of $Q$. Then $0\notin X$
and $X\cap S=\emptyset$. We claim that there are $x\in X$ and $s\in S$ with $x-s\notin X$. If not, then $X-S\subseteq
X$, that is, $X+S\subseteq X$ as $S=-S$; adding elements of $S$ repeatedly gives $X+\langle S\rangle\subseteq X$, and as $S$ generates $A$ and $X$ is not empty,
$X=A$, contradicting $0\notin X$. Let $t=s-x$. Then $t\ne 0$ because $x\notin S$. A new value $x'+t$ with $x'\in X$ is $0$
only if $x'=x-s$, which is not in $X$. So $f'=f+t\chi_Q$ is a flow. An edge of $Q$ with value $x$ has the new
value $s$ and is small in $f'$; the edges outside $Q$ did not change; and no edge of $Q$ was small in $f$. So $\rho(f')\le\rho(f)-1$.
\end{proof}

\begin{corollary}\label{cor:forest}
Every flow $f$ of a graph $G$ is joined by a chain of at most $\rho(f)$ steps to a flow in which the edges that are
not small form a forest.
\end{corollary}
\begin{proof}
Apply Lemma~\ref{lem:Astar} as long as the edges that are not small contain a cycle. Each application is a single step and lowers $\rho$ by at least $1$, and $\rho\ge 0$.
\end{proof}

\begin{lemma}[a cycle of small edges]\label{lem:Bstar}
Let $G$ be a loopless graph in which every vertex has degree $0$ or $4$. Let $f$ be a flow of $G$ with no balanced
pair, and $Q$ a cycle of $G$ all of whose edges are small. Then there are a direction of $Q$ and an element $t\ne 0$ of $A$ such that $f+t\chi_Q$ is a flow with a balanced pair.
\end{lemma}
\begin{proof}
Let $v$ be a vertex of $Q$, to be chosen below. As $v$ lies on $Q$, its degree is $4$; number the four ends at $v$ as $1,2,3,4$, where $1$ and $2$ are the ends at $v$ of the two edges of $Q$ at $v$, and write
$y_i=y_i(f)$. Direct $Q$ so that it enters $v$ through the end $1$ and leaves through the end $2$. For $j\in\{3,4\}$
let $t_j=y_1+y_j$, and let $j'$ be the other element of $\{3,4\}$; $t_j$ is not $0$ because $f$ has no balanced pair. By \eqref{eq:ends} and
\eqref{eq:kirchhoff}, in $f+t_j\chi_Q$ the end $1$ has the outward value $y_1-t_j=-y_j$ and the end $2$ has the
outward value $y_2+t_j=-y_{j'}$. Neither is $0$, and the ends $1$ and $j$ are a balanced pair. An edge of $Q$ that is not at $v$ gets
the value $0$ exactly if its value in the direction of $Q$ was $-t_j$. Hence the following holds.
\begin{itemize}
\item[$(\ast)$] If no edge of $Q$ that is not at $v$ has the value $-t_j$ in the direction of $Q$, then
$f+t_j\chi_Q$ is a flow with a balanced pair.
\end{itemize}
In particular, $(\ast)$ applies if $t_j$ is not small, because then $-t_j$ is not small either ($S=-S$), while all
values on $Q$ are small; and $(\ast)$ applies if $Q$ has only two edges, because then both edges of $Q$ are at $v$. Below
we use the identity $t_3+t_4=2y_1+y_3+y_4=y_1-y_2$, which holds by \eqref{eq:kirchhoff}.

\emph{Case 1: two consecutive edges of $Q$ have different colors.} Choose $v$ to be a common vertex of these two edges. Then $y_1=a$ is
small, of color $h$, and $y_2=b$ is small, of color $h'\ne h$. Suppose that $t_3$ and $t_4$ are both small. Their sum
$a-b$ has exactly two non-zero parts, $a$ in $\langle s_h\rangle$ and $-b$ in $\langle s_{h'}\rangle$. Two small elements
of the same color have a sum with at most one non-zero part; so $t_3$ and $t_4$ have different colors; by the uniqueness of the parts, $t_3$ and $t_4$ are the two non-zero parts of $t_3+t_4=a-b$, that is, $\{t_3,t_4\}=\{a,-b\}$. But $t_3=a$ means $y_3=0$, and $t_4=a$ means $y_4=0$. This is
impossible. So $t_3$ or $t_4$ is not small, and $(\ast)$ gives the step.

\emph{Case 2: all edges of $Q$ have the same color $h$.} Here $v$ may be any vertex of $Q$. Then $y_1,y_2\in\{s_h,-s_h\}$
and $y_1+y_2\ne 0$, so $y_1=y_2=a$ and $2a\ne 0$. (If $2s_h=0$ this case does not occur: the ends $1,2$ would be a
balanced pair.) The values of the edges of $Q$ in the direction of $Q$ are $a$ or $-a$. Now $-t_3=a$ means
$y_3=-2a$ and then $y_4=-2a-y_3=0$; and $-t_3=-a$ means $y_3=0$. Both are impossible, so $(\ast)$ applies with $j=3$ and gives a step with value $t_3$.
\end{proof}

\begin{proposition}[reaching a balanced pair when all degrees are $0$ or $4$]\label{prop:local4}
Let $G$ be a loopless graph in which every vertex has degree $0$ or $4$, and which has at least one edge. Every flow $f$ of $G$ is joined by a chain of at most $\rho(f)+1$ steps to a flow with a balanced pair.
\end{proposition}
\begin{proof}
By Corollary~\ref{cor:forest}, $f$ is joined by a chain of at most $\rho(f)$ steps to a flow $f^*$ in which the
edges that are not small form a forest. Let $n\ge 1$ be the number of vertices of degree $4$. Then $G$ has $2n$ edges, all on these
$n$ vertices; at most $n-1$ of them are not small in $f^*$ (a forest on $n$ vertices has at most $n-1$ edges), so at
least $n+1$ are small, and the small edges contain a cycle $Q$. If $f^*$ has no balanced pair,
Lemma~\ref{lem:Bstar} gives a further step to a flow with a balanced pair.
\end{proof}

\section{Finite abelian groups: proofs of Theorems \ref{thm:deg4} and \ref{thm:all}}\label{sec:proofs}

Theorem~\ref{thm:deg4} follows from Sections~\ref{sec:red} and~\ref{sec:local}. For Theorem~\ref{thm:all}, after the loops are
deleted (Lemma~\ref{lem:loop}), a vertex of degree at least $6$ is replaced by a triangle; Lemma~\ref{lem:triangle} shows that if the new graph is
$A$-flow-connected, then so is the original graph, and induction on $\sum_v\max(\deg v-4,0)$, the total amount by which the degrees exceed $4$, completes the proof.

\begin{proof}[Proof of Theorem~\ref{thm:deg4}]
If $|A|=1$, then $G$ has at most one flow (none if it has an edge), so $G$ is $A$-flow-connected. Let $|A|\ge 2$, let $\cG$ be the class of even graphs with all degrees at most $4$, and let $T=A\setminus\{0\}$; it satisfies \textup{(T1)} and \textup{(T2)}. Deleting a loop lowers the degree of its vertex by $2$. Splitting off as in Lemma~\ref{lem:split} lowers the degree of $v$ by $2$ and keeps the other degrees
(a new loop at $u=w$ counts twice). So $\cG$ is closed under both operations. A graph in $\cG$ with no loop and no
vertex of degree $2$ has only vertices of degree $0$ or $4$, so \textup{(LOCAL)} is Proposition~\ref{prop:local4}.
Proposition~\ref{prop:red} gives the theorem.
\end{proof}

\begin{lemma}[contracting a triangle]\label{lem:triangle}
Let $|A|\ge 4$. Let $G^+$ be a graph and $\Delta$ a triangle of $G^+$ with vertices $x_1,x_2,x_3$, such that the
three edges of $\Delta$ are the only edges of $G^+$ whose two ends are both at vertices of $\Delta$. Let $G=G^+/\Delta$ be the
graph obtained by deleting the three edges of $\Delta$ and identifying $x_1,x_2,x_3$ to one vertex $x$. For a flow
$f$ of $G^+$, let $pf$ be its restriction to $E(G)=E(G^+)\setminus E(\Delta)$.
\begin{itemize}
\item[(a)] $pf$ is a flow of $G$, and every flow of $G$ is $pf$ for some flow $f$ of $G^+$.
\item[(b)] If $f_2$ comes from $f_1$ by a multi-step with value $t$, then $pf_2=pf_1$ or $pf_2$ comes from $pf_1$
by a multi-step with value $t$.
\end{itemize}
So if $G^+$ is $A$-flow-connected, then $G$ is $A$-flow-connected.
\end{lemma}
\begin{proof}
(a) The flow condition of $pf$ at $x$ is the sum of the flow conditions of $f$ at $x_1,x_2,x_3$, in which the edges
of $\Delta$ cancel. Let $\varphi$ be a flow of $G$. For $i=1,2,3$, let $\beta_i$ be the sum of the outward values of $\varphi$
over the ends at $x$ that are ends at $x_i$ in $G^+$; then $\beta_1+\beta_2+\beta_3=0$. Give the edges $x_1x_2$,
$x_2x_3$, $x_3x_1$ of $\Delta$, in these directions, the values $z$, $z-\beta_2$, $z+\beta_1$. The flow conditions at
$x_1,x_2,x_3$ are $\beta_1+z-(z+\beta_1)=0$, $\beta_2+(z-\beta_2)-z=0$ and $\beta_3+(z+\beta_1)-(z-\beta_2)=0$; they
hold for every $z$. The three values are not $0$ if $z\notin\{0,\beta_2,-\beta_1\}$, and $|A|\ge 4$ leaves such a
$z$.

(b) Let $f_2=f_1+t(\chi_{Q_1}+\dots+\chi_{Q_r})$ with directed cycles of $G^+$, no two of which have a common edge.
Let $H$ be the directed graph on $V(G)$ whose arcs are the edges of $Q_1,\dots,Q_r$ outside $\Delta$, each with the
direction of its cycle. At every vertex other than $x$, the in-degree of $H$ equals its out-degree, because this holds for the union of the
cycles. It also holds at $x$: in the union of the cycles, the total in-degree of $x_1,x_2,x_3$ equals their total
out-degree, and every arc of the cycles inside $\Delta$ adds one to the total in-degree and one to the total out-degree. So $H$ is the union of
directed cycles of $G$, no two of which have a common arc, and $pf_2-pf_1$ is $t$ times the sum of their maps $\chi$.
If $H$ has no arc, $pf_2=pf_1$; otherwise this is a multi-step of $G$ with value $t$.

The last statement follows, because a multi-step is a chain of steps and $p$ is onto by (a).
\end{proof}

\begin{proof}[Proof of Theorem~\ref{thm:all}]
By Lemma~\ref{lem:loop} (with $T=A\setminus\{0\}$), deleting the loops one by one, we may assume that $G$ has no loop. Let
$\Phi(G)=\sum_v\max(\deg v-4,0)$. We use induction on $\Phi(G)$. If $\Phi(G)=0$, this is Theorem~\ref{thm:deg4}.
Otherwise let $v$ be a vertex of degree $d\ge 6$. Let $G^+$ be the graph obtained from $G$ by replacing $v$ with a
triangle $\Delta$ with vertices $x_1,x_2,x_3$: two of the ends at $v$ become ends at $x_1$, two become ends at $x_2$,
and the other $d-4$ become ends at $x_3$. The new degrees are $4$, $4$ and $d-2$, so $G^+$ is even and
$\Phi(G^+)=\Phi(G)-2$. The graph $G^+$ has no loop, and the edges of $\Delta$ are the only edges whose two ends
are both at vertices of $\Delta$, because $G$ has no loop at $v$. Also $G=G^+/\Delta$. By the induction hypothesis $G^+$ is
$A$-flow-connected, and Lemma~\ref{lem:triangle} gives the same for $G$.
\end{proof}

\section{Cyclic groups: steps with values \texorpdfstring{$\pm1$}{+-1} and \texorpdfstring{$\pm2$}{+-2}}\label{sec:cyclic}

This section proves Theorem~\ref{thm:cyclic}. Let $A=\Z_m$ with $m\ge 4$ and $T=\{1,-1,2,-2\}$. An edge is small for a flow $f$ if its value is $1$
or $-1$ (as in Section~\ref{sec:local} with $S=\{1,-1\}$), and $\rho(f)$ is the number of edges that are not small. The value of an edge that is not
small is, in each direction, the residue of an integer in $\{2,\dots,m-2\}$.

\begin{lemma}[a cycle with no small edge, in $\Z_m$]\label{lem:Acyc}
Let $f$ be a flow of a graph $G$ and $Q$ a cycle none of whose edges is small. Then $f$ is $T$-joined, by steps along
$Q$ with value $-1$, to a flow $f'$ with $\rho(f')\le\rho(f)-1$.
\end{lemma}
\begin{proof}
Fix a direction of $Q$. For an edge $e$ of $Q$, let $x_e\in\{2,\dots,m-2\}$ be the integer whose residue is the value
of $e$ in this direction, and let $x$ be the least of the $x_e$. For $j=1,\dots,x-1$ the map $f-j\chi_Q$ has on $e$
the residue of $x_e-j$, where $1\le x_e-j\le m-3$; so it is a flow. In $f'=f-(x-1)\chi_Q$, an edge with $x_e=x$
has the value $1$; the edges outside $Q$ did not change, and no edge of $Q$ was small in $f$. So $\rho(f')\le\rho(f)-1$.
\end{proof}

\begin{lemma}[a cycle of small edges, in $\Z_m$]\label{lem:Bcyc}
Let $G$ have no loop, let $f$ be a flow with no balanced pair, $Q$ a cycle all of whose edges are small, and $v$ a
vertex of $Q$.
\begin{itemize}
\item[(a)] The two ends of $Q$ at $v$ have the same outward value $a\in\{1,-1\}$.
\item[(b)] If an end at $v$ that is not an end of $Q$ has an outward value $c\ne -2a$, then $f$ is $T$-joined, by steps along $Q$, to a
flow with a balanced pair at $v$.
\end{itemize}
\end{lemma}
\begin{proof}
(a) Both values are in $\{1,-1\}$ and their sum is not $0$.

(b) Call the two ends of $Q$ at $v$ the end $1$ and the end $2$, direct $Q$ so that it enters $v$ through the end
$1$ and leaves through the end $2$, and let $t=a+c$. Then $t\ne
0$ ($f$ has no balanced pair), $t\ne a$ ($c\ne 0$), and $t\ne -a$ ($c\ne -2a$). So $t$ is the residue of an integer
$\tau$ with $2\le\tau\le m-2$. The values of the edges of $Q$ in the direction of $Q$ are $1$ or $-1$. For
$j=2,3,\dots,\tau$ the map $f+j\chi_Q$ has on the edges of $Q$ the residues of $1+j$ or $-1+j$, integers between $1$ and
$m-1$; so it is a flow. The chain $f$, $f+2\chi_Q$, $f+3\chi_Q$, \dots, $f+\tau\chi_Q$ consists of one step with
value $2$ and $\tau-2$ steps with value $1$. In $f+\tau\chi_Q=f+t\chi_Q$ the end $1$ has the outward value $a-t=-c$ by
\eqref{eq:ends}, and the end with the outward value $c$ did not change; so these two ends are a balanced pair at $v$.
\end{proof}

\begin{proposition}[reaching a balanced pair in $\Z_m$]\label{prop:localcyc}
Let $m\ge 4$, and let $G$ be a loopless even graph in which every vertex has degree $0$ or at least $4$, and which has at least one edge. Every flow of $G$ is $T$-joined to a flow with a balanced pair.
\end{proposition}
\begin{proof}
Let $n\ge 1$ be the number of vertices of degree at least $4$; all edges are on these vertices, and $|E(G)|\ge 2n$.
By Lemma~\ref{lem:Acyc}, used as long as the edges that are not small contain a cycle, a given flow is $T$-joined to a flow
$f^*$ in which they form a forest; so at most $n-1$ edges are not small, as they lie on the $n$ vertices. Suppose that
$f^*$ has no balanced pair (otherwise there is nothing to prove), and that the hypothesis of Lemma~\ref{lem:Bcyc}~(b) fails for every cycle $Q$ of small
edges and every vertex $v$ of $Q$. Then at every vertex $v$ of such a cycle, all ends at $v$ that are not ends of the cycle have the outward
value $-2a$, where $a\in\{1,-1\}$ is the common outward value at $v$ of the two ends of the cycle
(Lemma~\ref{lem:Bcyc}~(a)). As $m\ge 4$, $-2a$ is neither $1$ nor $-1$: $-2a=a$ would mean $3\equiv 0$, and $-2a=-a$ would
mean $1\equiv 0$, modulo $m$. So the edges of these ends are not small, and every vertex of a cycle of small edges lies
on exactly two small edges, the two edges of the cycle. Hence every cycle of small edges is a whole component of the graph formed by the small edges on the $n$
vertices. So every component of this graph is a tree or a cycle (a component that contains a cycle is that cycle), and
there are at most $n$ small edges. Then
$|E(G)|\le n+(n-1)<2n$, a contradiction. So, for some cycle $Q$ of small edges and some vertex $v$ of $Q$,
Lemma~\ref{lem:Bcyc}~(b) joins $f^*$ to a flow with a balanced pair.
\end{proof}

\begin{proof}[Proof of Theorem~\ref{thm:cyclic}]
Apply Proposition~\ref{prop:red} to the class $\cG$ of all even graphs and $T=\{1,-1,2,-2\}$: \textup{(T1)} and
\textup{(T2)} hold, as noted at the start of Section~\ref{sec:red}, and \textup{(LOCAL)} is
Proposition~\ref{prop:localcyc}.
\end{proof}

\begin{remark}\label{rem:noncyclic}
Theorem~\ref{thm:cyclic} proves Theorem~\ref{thm:all} again for cyclic groups, without Lemma~\ref{lem:triangle}. The
counting argument of Proposition~\ref{prop:localcyc} does not extend to all groups:
it requires the value $-2a$ in Lemma~\ref{lem:Bcyc}~(b) not to be small. This holds in $\Z_m$ with $m\ge 4$, but fails,
for example, in $\Z_3\times\Z_3$ with $S$ as in Section~\ref{sec:local}, where $-2a=a$; there a vertex of degree $6$ can have six ends
with the same small outward value $a$ and no balanced pair. For this reason, the proof of Theorem~\ref{thm:all} first treats graphs with all degrees at most $4$
(Section~\ref{sec:local} and Theorem~\ref{thm:deg4}) and then reduces larger degrees by triangles
(Section~\ref{sec:proofs}).
\end{remark}

\section{Integer flows: proof of Theorem~\ref{thm:main}~(b)}\label{sec:int}

This section turns a chain of $\Z_k$-flows from Theorem~\ref{thm:cyclic}
into a walk in $\cF(G,k)$: every $\Z_k$-flow of the chain is replaced by a path of $k$-flows (integer flows with values between $1$ and $k-1$ in absolute value, defined below) whose residues modulo $k$
form that $\Z_k$-flow.

\paragraph{Definitions.} An \emph{integer flow} of a graph $G$ is a map $F\colon E(G)\to\Z$ that satisfies the flow condition at every vertex. For a map $x$ from $E(G)$ to $\Z$ or to the real numbers, the \emph{boundary}
$(\partial x)(v)$ at a vertex $v$ is the sum of $x$ over the edges that enter $v$ minus the sum over the edges that
leave $v$ (a loop at $v$ is in both sums). So an integer map $x$ is an integer flow if and only if $\partial x=0$; also $\partial\chi_Q=0$ for every directed cycle~$Q$.

Let $k\ge 2$. A \emph{$k$-flow} of $G$ is an integer flow $F$ with $1\le|F(e)|\le k-1$ for every edge $e$; for $k=3$, these are the nowhere-zero $3$-flows of Section~\ref{sec:intro}. The graph
$\cF(G,k)$ has the $k$-flows of $G$ as its vertices, and two $k$-flows are adjacent if they differ by $\lambda\chi_Q$,
where $Q$ is a directed cycle and $\lambda\ne 0$ is an integer. The graph $G$ is \emph{$k$-flow-connected} if $\cF(G,k)$ is
connected; as for $\cF(G,A)$, it counts as connected when it has at most one vertex. As with $\cF(G,A)$ in Section~\ref{sec:words}, this graph is the same as the graph $\cF(G,k)$ of
\cite{CLSWX,EHLMNS}. In this section a $\Z_k$-flow means a nowhere-zero $\Z_k$-flow. An \emph{integer lift} of a $\Z_k$-flow $f$ is a $k$-flow $F$ with $F(e)\equiv f(e)$
modulo $k$ for every edge $e$. The residues modulo $k$ of a $k$-flow form a $\Z_k$-flow, because no value of a
$k$-flow is a multiple of $k$.

\begin{theorem}\label{thm:int}
Let $k\ge 4$. Every even graph is $k$-flow-connected.
\end{theorem}

This is part~(b) of Theorem~\ref{thm:main}. The proof lifts a chain from Theorem~\ref{thm:cyclic} step by step.
Lemma~\ref{lem:lift} is Tutte's lifting theorem in the form of \cite[Theorem~2.3]{EHLMNS}, and Lemma~\ref{lem:fibre}
is contained in \cite[Lemma~5.5]{EHLMNS}. We include their short proofs to keep the paper self-contained; the rounding lemma
(Lemma~\ref{lem:round}) gives Lemma~\ref{lem:lift} and is used again in the proof of Lemma~\ref{lem:step2}.

\begin{lemma}[rounding]\label{lem:round}
Let $G$ be a graph and $x$ a map from $E(G)$ to the real interval $[0,1]$ such that $(\partial x)(v)$ is an integer
for every vertex $v$. Then there is a map $z\colon E(G)\to\{0,1\}$ with $\partial z=\partial x$ and with $z(e)=x(e)$
for every edge $e$ with $x(e)\in\{0,1\}$.
\end{lemma}
\begin{proof}
Call an edge $e$ \emph{fractional} if $0<x(e)<1$. If no edge is fractional, take $z=x$. Otherwise the fractional
edges contain a cycle $Q$. If some loop is fractional, it is such a cycle. If no loop is fractional, then no vertex $v$
meets exactly one fractional edge, because the other edges at $v$ have integer values and $(\partial x)(v)$ is an
integer; and a graph with no loop, at least one edge and no vertex of degree $1$ contains a cycle. Give $Q$ a direction, and
let $\delta>0$ be the least number such that $x+\delta\chi_Q$ has the value $0$ or $1$ on some edge of $Q$. The map
$x+\delta\chi_Q$ takes values in $[0,1]$, has the same boundary as $x$, agrees with $x$ on the edges that are not
fractional, and has fewer fractional edges than $x$. After at most $|E(G)|$ such changes no edge is fractional; take
$z$ to be the final map.
\end{proof}

For a $\Z_k$-flow $f$ and an edge $e$, let $r_f(e)\in\{1,\dots,k-1\}$ be the integer whose residue is the value of
$f$ on $e$ in the reference direction. Then $\partial r_f\equiv 0$ modulo $k$ at every vertex, and an integer lift of
$f$ has on every edge $e$ the value $r_f(e)$ or $r_f(e)-k$.

\begin{lemma}[integer lifts]\label{lem:lift}
Let $k\ge 2$. Every $\Z_k$-flow $f$ of a graph $G$ has an integer lift.
\end{lemma}
\begin{proof}
The map $x=r_f/k$ takes values in $(0,1)$ and has integer boundary. Let $z$ be as in Lemma~\ref{lem:round}, and
$F=r_f-kz$. Then $\partial F=\partial r_f-k\,\partial z=\partial r_f-k\,\partial x=0$, $F\equiv f$ modulo $k$, and
$F(e)$ is $r_f(e)\in[1,k-1]$ or $r_f(e)-k\in[-(k-1),-1]$.
\end{proof}

\begin{lemma}[integer lifts of a given flow]\label{lem:fibre}
Let $k\ge 2$ and let $f$ be a $\Z_k$-flow of a graph $G$. Any two integer lifts of $f$ are joined by a path in
$\cF(G,k)$ all of whose vertices are integer lifts of $f$.
\end{lemma}
\begin{proof}
Let $F,F'$ be integer lifts of $f$. On every edge both have the value $r_f(e)$ or $r_f(e)-k$, so $h=(F'-F)/k$ is an
integer flow with values in $\{-1,0,1\}$. Direct every edge $e$ with $h(e)\ne 0$ so that $h$ is $1$ in this
direction. By the flow condition these arcs form a directed graph with in-degree equal to out-degree at every vertex,
so they are the union of directed cycles $Q_1,\dots,Q_j$, no two of which have a common arc, and
$h=\chi_{Q_1}+\dots+\chi_{Q_j}$. Let $F_i=F+k(\chi_{Q_1}+\dots+\chi_{Q_i})$. On every edge $F_i$ has the value of $F$
or the value of $F'$, so $F_i$ is an integer lift of $f$; and $F_i-F_{i-1}=k\chi_{Q_i}$. So $F=F_0,F_1,\dots,F_j=F'$
is a path in $\cF(G,k)$ of integer lifts of $f$.
\end{proof}

\begin{lemma}[a step with value $1$]\label{lem:step1}
Let $k\ge 2$, let $Q$ be a directed cycle of a graph $G$, and let $f$ and $f+\chi_Q$ be $\Z_k$-flows of $G$. Then for
every integer lift $F$ of $f$, the map $F+\chi_Q$ is an integer lift of $f+\chi_Q$.
\end{lemma}
\begin{proof}
Choose the reference directions so that $Q$ traverses all its edges in the reference direction. (The statement does
not depend on this choice: a change of the reference direction of an edge changes the signs of $F(e)$, of $f(e)$ and
of $\chi_Q(e)$.) The map $F+\chi_Q$ is an integer flow with the residues of $f+\chi_Q$. Let $e$ be an edge of $Q$. Then
$r_f(e)\ne k-1$, because $f+\chi_Q$ is not $0$ on $e$. The value $F(e)$ is $r_f(e)$ or $r_f(e)-k$, so $F(e)+1$ is in
$[2,k-1]$ or in $[2-k,-1]$.
\end{proof}

\begin{lemma}[a step with value $2$]\label{lem:step2}
Let $k\ge 3$, let $Q$ be a directed cycle of a graph $G$, and let $f$ and $f+2\chi_Q$ be $\Z_k$-flows of $G$. Then
there is an integer lift $F$ of $f$ such that $F+2\chi_Q$ is an integer lift of $f+2\chi_Q$.
\end{lemma}
\begin{proof}
Lemma~\ref{lem:step1} cannot simply be used twice, because $f+\chi_Q$ can have zeros. As in the proof of Lemma~\ref{lem:step1}, choose the reference directions so that $Q$ traverses all its edges in the reference
direction, and write $r=r_f$. No edge of $Q$ has $r(e)=k-2$, because $f+2\chi_Q$ is not $0$ there. Let $E_1$ be the set of edges $e$ of $Q$ with $r(e)=k-1$, and let
\[
x=(r+\chi_Q)/k .
\]
All values of $x$ are in $(0,1]$, and $x(e)=1$ exactly for $e\in E_1$. The boundary $\partial x=(\partial r)/k$ is
an integer at every vertex. By Lemma~\ref{lem:round} there is $z\colon E(G)\to\{0,1\}$ with $\partial z=\partial x$
and with $z(e)=1$ for every $e\in E_1$. Let $F=r-kz$. As in the proof of Lemma~\ref{lem:lift}, $F$ is an integer
lift of $f$.

Now $F+2\chi_Q$ is an integer flow with the residues of $f+2\chi_Q$, and outside $Q$ it equals $F$. If $e\in E_1$, then $z(e)=1$, so
$F(e)=-1$ and the new value is $1$. If $e$ is an edge of $Q$ not in $E_1$, then $r(e)\le k-3$; hence
$3\le F(e)+2\le k-1$ if $F(e)=r(e)$, and $3-k\le F(e)+2\le -1$ if $F(e)=r(e)-k$. So $F+2\chi_Q$ is a $k$-flow.
\end{proof}

\begin{proposition}[from $\Z_k$-flows to $k$-flows]\label{prop:transfer}
Let $k\ge 3$ and let $G$ be a graph. If any two $\Z_k$-flows of $G$ are joined by a chain of steps whose values are in
$\{1,-1,2,-2\}$, then $G$ is $k$-flow-connected.
\end{proposition}
\begin{proof}
Let $F$ and $F'$ be $k$-flows of $G$, and let $f$ and $f'$ be the $\Z_k$-flows formed by their residues modulo $k$. If $f=f'$, then $F$
and $F'$ are two integer lifts of $f$, and they are joined in $\cF(G,k)$ by Lemma~\ref{lem:fibre}. Otherwise let $f=f_0,f_1,\dots,f_s=f'$
with $s\ge 1$ be a chain of $\Z_k$-flows with $f_i=f_{i-1}+t_i\chi_{Q_i}$ and $t_i\in\{1,-1,2,-2\}$. If $t_i$ is
neither $1$ nor $2$, replace $Q_i$ by the same cycle with the other direction and $t_i$ by $-t_i$; so we may assume
that $t_i\in\{1,2\}$. For every $i$ from $1$ to $s$, Lemmas~\ref{lem:lift} and~\ref{lem:step1} (if $t_i=1$) or Lemma~\ref{lem:step2} (if
$t_i=2$) give an integer lift $F_{i-1}^{+}$ of $f_{i-1}$ such that $F_i^{-}=F_{i-1}^{+}+t_i\chi_{Q_i}$ is an integer
lift of $f_i$; these two $k$-flows are adjacent in $\cF(G,k)$. (The superscripts $-$ and $+$ only label two, possibly equal, integer lifts of $f_i$: $F_i^{-}$ is the lift reached by the step from $f_{i-1}$, and $F_i^{+}$ is the lift from which the step to $f_{i+1}$ starts.) By Lemma~\ref{lem:fibre}, $F$ is
joined to $F_0^{+}$, $F_i^{-}$ is joined to $F_i^{+}$ for $0<i<s$, and $F_s^{-}$ is joined to $F'$. So $F$ and $F'$
are joined in $\cF(G,k)$.
\end{proof}

\begin{proof}[Proof of Theorem~\ref{thm:int}]
By Theorem~\ref{thm:cyclic} with $m=k$, any two $\Z_k$-flows of an even graph are joined by a chain of steps
with values in $\{1,-1,2,-2\}$; so Proposition~\ref{prop:transfer} shows that the graph is $k$-flow-connected.
\end{proof}


\section*{Acknowledgments}
This work was done overwhelmingly by AI models and supervised by the author. The models were Anthropic's Claude and
OpenAI's ChatGPT. The author takes full responsibility for the paper.

\begin{thebibliography}{9}\small
\bibitem{CLSWX} D.~W. Cranston, J.~Li, B.~Su, Z.~Wang, N.~Xu, \emph{Reconfiguration of nowhere-zero flows},
arXiv:2606.24685v1 (June 23, 2026). Statement numbers refer to this version.
\bibitem{EHLMNS} L.~Esperet, K.~Hendrey, A.~Lagoutte, M.~Marseloo, S.~Norin, R.~Steiner, \emph{Nowhere-zero flow
reconfiguration}, arXiv:2512.17342v4 (July 3, 2026). Statement numbers refer to this version.
\bibitem{G} D.~Goel, \emph{The Cranston--Li--Su--Wang--Xu (2026) conjecture is true, with extensions to graphs with
few vertices of odd degree}, Zenodo (2026),
\href{https://doi.org/10.5281/zenodo.23605838}{doi:10.5281/zenodo.23605838}. Statement numbers refer to version 2.30,
\href{https://doi.org/10.5281/zenodo.23605839}{doi:10.5281/zenodo.23605839}.
\bibitem{JLPT} F.~Jaeger, N.~Linial, C.~Payan, M.~Tarsi, \emph{Group connectivity of graphs---a nonhomogeneous
analogue of nowhere-zero flow properties}, J. Combin. Theory Ser.~B \textbf{56} (1992), no.~2, 165--182.
\bibitem{Tutte} W.~T.~Tutte, \emph{A contribution to the theory of chromatic polynomials}, Canad. J. Math.
\textbf{6} (1954), 80--91.
\end{thebibliography}
\end{document}
